\begin{corollary}[6.9.3]
\label{IV.2.s6.9.3-fr}
Sean $S$ un preesquema noetheriano, $u:X\to S$ un morfismo de tipo finito,
$\mathcal F$ un $\mathcal O_X$-Módulo coherente. Existe entonces una partición
$(S_i)_{1\leq i\leq n}$ de $S$ en un número finito de conjuntos localmente cerrados
en $S$, tal que, si se designa de nuevo por $S_i$ el subpreesquema reducido
de $S$ cuyo espacio subyacente es $S_i$, y si se pone
$X_i=X\times_S S_i$, el $\mathcal O_{X_i}$-Módulo
$\mathcal F_i=\mathcal F\otimes_{\mathcal O_S}\mathcal O_{S_i}$ sea plano
sobre $S_i$.
\end{corollary}

\begin{proof}
Procedamos por inducción noetheriana $(\mathbf{0}_{\mathrm I},2.2.2)$ sobre el conjunto
de las partes cerradas $T$ de $S$ tales que el subpreesquema reducido de $S$
cuyo espacio subyacente es $T$ verifica la conclusión de (6.9.3). Es posible,
pues, limitarse a demostrar el corolario para $S$ suponiendo que es cierto
para todo subpreesquema cerrado reducido de $S$ cuyo espacio subyacente es una
parte cerrada $\ne S$. Como el morfismo $S_{\mathrm{red}}\to S$ es de tipo
finito y sobreyectivo, se puede reemplazar $S$ por $S_{\mathrm{red}}$ y suponer,
por tanto, que $S$ es reducido y no vacío. Como $S$ es noetheriano, el interior $T$ de una
componente irreducible de $S$ no es vacío, y el preesquema inducido sobre $T$
es íntegro; en virtud de (6.9.1), existe entonces un abierto no vacío $U\subset T$
\oldpage[IV-2]{155}
tal que $\mathcal F|u^{-1}(U)$ es plano sobre $U$. Si $Y$ es el
subpreesquema reducido de $S$ cuyo espacio subyacente es $S\setmin U$, por
hipótesis existe una partición $(Y_i)$ de $Y$ en conjuntos localmente cerrados en
$Y$ (y, por tanto, en $S$) tales que
$\mathcal F_i=\mathcal F\otimes_{\mathcal O_S}\mathcal O_{Y_i}$ es plano
sobre $Y_i$ para todo $i$; es claro que los $Y_i$ y $U$ forman la
partición buscada.
\end{proof}
